\documentclass[11pt,a4paper]{article}
\usepackage{notebooks}

\notebooktitle{Linux Commands}
\notebooksubtitle{Permissions, Users and Groups}
\notebooktagline{Why "Permission denied" happens, and what to do about it.}
\notebookauthor{Bhuvnesh Verma}
\notebooksource{github.com/MasterBhuvnesh/notebooks}

\begin{document}
\notebookcover

\section{File Permissions}

In Linux, \textbf{permissions (modes)} control who can \textbf{read, write, or execute} a
file or directory.

When you use:

\begin{sh}
mkdir -m 755 myapp
\end{sh}

the \lstinline|755| is the \textbf{permission mode} assigned to the directory.

\subsection{Permission Types}

There are three permission bits:

\begin{center}
\small
\begin{tabular}{@{}l l l@{}}
\toprule
\rowcolor{tablehead}\bfseries Permission & \bfseries Symbol & \bfseries Value \\
\midrule
Read    & \texttt{r} & 4 \\
Write   & \texttt{w} & 2 \\
Execute & \texttt{x} & 1 \\
\bottomrule
\end{tabular}
\end{center}

The values are added together:

\begin{center}
\small
\begin{tabular}{@{}l l l@{}}
\toprule
\rowcolor{tablehead}\bfseries Permission Set & \bfseries Calculation & \bfseries Value \\
\midrule
\texttt{-{}-{}-} & 0     & 0 \\
\texttt{-{}-x}   & 1     & 1 \\
\texttt{-w-}     & 2     & 2 \\
\texttt{-wx}     & 2+1   & 3 \\
\texttt{r-{}-}   & 4     & 4 \\
\texttt{r-x}     & 4+1   & 5 \\
\texttt{rw-}     & 4+2   & 6 \\
\texttt{rwx}     & 4+2+1 & 7 \\
\bottomrule
\end{tabular}
\end{center}

\subsection{Who Gets Permissions?}

Linux divides users into three groups:

\begin{center}
\small
\begin{tabular}{@{}l l@{}}
\toprule
\rowcolor{tablehead}\bfseries Group & \bfseries Meaning \\
\midrule
User (\texttt{u})   & Owner of the file \\
Group (\texttt{g})  & Users in the file's group \\
Others (\texttt{o}) & Everyone else \\
\bottomrule
\end{tabular}
\end{center}

A permission like \lstinline|755| means:

\begin{lstlisting}[language=]
7   5   5
│   │   │
│   │   └── Others
│   └────── Group
└────────── Owner
\end{lstlisting}

\subsection{Example: 755}

Break it down:

\begin{lstlisting}[language=]
7 = rwx = 4+2+1
5 = r-x = 4+1
5 = r-x = 4+1
\end{lstlisting}

Result: \lstinline|rwxr-xr-x|. Meaning:

\begin{center}
\small
\begin{tabular}{@{}l l@{}}
\toprule
\rowcolor{tablehead}\bfseries User Type & \bfseries Permissions \\
\midrule
Owner  & Read, Write, Execute \\
Group  & Read, Execute \\
Others & Read, Execute \\
\bottomrule
\end{tabular}
\end{center}

\subsection{Example: 700}

\begin{lstlisting}[language=]
7 = rwx
0 = ---
0 = ---
\end{lstlisting}

Result: \lstinline|rwx------|. Only the owner can access it. Useful for:

\begin{sh}
mkdir -m 700 secrets
\end{sh}

\subsection{Example: 777}

Result: \lstinline|rwxrwxrwx|. Everyone can read, write and execute.

\field{Warning:}{Usually avoided, because anyone can modify the contents.}

\subsection{Example: 644}

Common for files.

\begin{lstlisting}[language=]
6 = rw-
4 = r--
4 = r--
\end{lstlisting}

Result: \lstinline|rw-r--r--|. The owner can read and write; everyone else can only read.

\subsection{What Does Execute Mean for a Directory?}

For a \textbf{file}, \texttt{x} means the file can be run as a program or script. For a
\textbf{directory}, \texttt{x} means you can enter it, that is \lstinline|cd| into it.

\begin{sh}
chmod 600 mydir
\end{sh}

That gives \lstinline|rw-------|. You can see the directory, but you \textbf{cannot}
\lstinline|cd| into it, because execute permission is missing. To enter a directory,
\texttt{x} permission is required.

\subsection{Viewing Permissions}

\begin{sh}
ls -l
\end{sh}

Example output:

\begin{lstlisting}[language=]
drwxr-xr-x 2 bhuvnesh users 4096 Aug 15 project
\end{lstlisting}

Breakdown:

\begin{lstlisting}[language=]
d rwx r-x r-x
│ │   │   │
│ │   │   └── Others
│ │   └────── Group
│ └────────── Owner
└──────────── Directory
\end{lstlisting}

\subsection{Changing Permissions}

Numeric method:

\begin{sh}
chmod 755 project
chmod 700 secrets
chmod 644 file.txt
\end{sh}

Symbolic method:

\begin{sh}
chmod u+x script.sh   # add execute permission for the owner
chmod g+w file.txt    # add write permission for the group
chmod o-r file.txt    # remove read permission for others
\end{sh}

\subsection{Common Modes You Will Use}

\begin{center}
\small
\begin{tabular}{@{}l l l@{}}
\toprule
\rowcolor{tablehead}\bfseries Mode & \bfseries Meaning & \bfseries Typical Use \\
\midrule
755 & \texttt{rwxr-xr-x} & Application directories \\
700 & \texttt{rwx-{}-{}-{}-{}-{}-} & Private directories \\
644 & \texttt{rw-r-{}-r-{}-} & Regular files \\
600 & \texttt{rw-{}-{}-{}-{}-{}-{}-} & Secrets, keys, passwords \\
777 & \texttt{rwxrwxrwx} & Rarely recommended \\
\bottomrule
\end{tabular}
\end{center}

\subsubsection*{Real-world examples}

\begin{sh}
mkdir -m 755 backend
mkdir -m 700 .ssh
touch app.js
chmod 644 app.js
chmod 600 id_rsa
\end{sh}

For software development you will most often see:

\begin{itemize}\itemsep2pt
  \item \textbf{Directories:} \lstinline|755|
  \item \textbf{Source code files:} \lstinline|644|
  \item \textbf{Shell scripts:} \lstinline|755|
  \item \textbf{SSH private keys and \lstinline|.env| files:} \lstinline|600| or \lstinline|400|
\end{itemize}

\newpage
\section{Users, Groups and Ownership}

To really understand Linux permissions, you need to understand \textbf{users},
\textbf{groups} and \textbf{ownership} first. Linux is a multi-user operating system. On a
company server:

\begin{lstlisting}[language=]
ubuntu
├── bhuvnesh
├── rohan
├── alice
└── john
\end{lstlisting}

Each person has their own account.

\subsection{Viewing the Current User}

\begin{sh}
whoami
# bhuvnesh

id
# uid=1000(bhuvnesh) gid=1000(bhuvnesh) groups=1000(bhuvnesh),27(sudo)
\end{sh}

\begin{itemize}\itemsep2pt
  \item UID = User ID
  \item GID = Group ID
  \item Groups = the groups this user belongs to
\end{itemize}

\subsection{What is a User?}

A user is an account that can log in, own files, run processes and hold permissions.

\begin{sh}
sudo useradd john    # create a user
sudo passwd john     # set the password
\end{sh}

\subsection{What is a Group?}

A group is a collection of users.

\begin{lstlisting}[language=]
developers
├── bhuvnesh
├── rohan
└── alice
\end{lstlisting}

Instead of assigning permissions to each user, assign them to the group.

\subsection{Managing Groups}

\begin{sh}
sudo groupadd developers          # create a group
cat /etc/group                    # check it exists

sudo usermod -aG developers john  # add a user (-a append, -G group)
groups john
# john : john developers

sudo gpasswd -d john developers   # remove a user from the group
sudo groupdel developers          # delete the group
\end{sh}

\subsection{Deleting a User}

\begin{sh}
sudo userdel john      # delete the user
sudo userdel -r john   # delete the user and their home directory
\end{sh}

\subsection{Where Users Are Stored}

Linux stores users in \lstinline|/etc/passwd|:

\begin{sh}
cat /etc/passwd
\end{sh}

\begin{lstlisting}[language=]
john:x:1001:1001::/home/john:/bin/bash
\end{lstlisting}

The fields are: username, password placeholder, UID, GID, comment, home directory, shell.

\subsection{Where Passwords Are Stored}

\begin{sh}
sudo cat /etc/shadow
\end{sh}

Only root can read it. It contains hashed passwords.

\subsection{What is Root?}

Linux has a superuser called \texttt{root}. Root can create and delete users, modify any
file, install software and shut the server down.

\begin{sh}
sudo su
# or
su -

whoami
# root
\end{sh}

\subsection{Ownership of Files}

\begin{sh}
touch app.js
ls -l
# -rw-r--r-- 1 bhuvnesh developers 0 Aug 15 app.js
\end{sh}

Here the owner is \texttt{bhuvnesh} and the group is \texttt{developers}.

\begin{sh}
sudo chown john app.js              # change the owner
sudo chgrp developers app.js        # change the group
sudo chown john:developers app.js   # change both at once
\end{sh}

\subsection{Directory Permissions Example}

\begin{sh}
mkdir project
chmod 770 project
\end{sh}

With owner \texttt{bhuvnesh} and group \texttt{developers}, \lstinline|770| means
\lstinline|rwxrwx---|:

\begin{center}
\small
\begin{tabular}{@{}l l@{}}
\toprule
\rowcolor{tablehead}\bfseries Who & \bfseries Access \\
\midrule
Owner  & Full \\
Group  & Full \\
Others & None \\
\bottomrule
\end{tabular}
\end{center}

So anyone in \texttt{developers} can use the directory.

\subsection{Real Company Example}

\begin{lstlisting}[language=]
developers          designers
├── bhuvnesh        ├── john
├── rohan           └── sara
└── alice
\end{lstlisting}

Create a shared directory:

\begin{sh}
sudo mkdir /shared-code
sudo chown root:developers /shared-code
sudo chmod 770 /shared-code
\end{sh}

Developers can access it; designers cannot.

\subsection{Sudo Access}

Not every user can become root. On Ubuntu, users allowed to run \lstinline|sudo| belong to
the \texttt{sudo} group.

\begin{sh}
groups bhuvnesh
# bhuvnesh sudo

sudo usermod -aG sudo john   # now john can run: sudo apt update
\end{sh}

\subsection{System Users vs Human Users}

Listing \lstinline|/etc/passwd| shows entries like \texttt{root}, \texttt{daemon},
\texttt{www-data}, \texttt{mysql}, \texttt{postgres}, \texttt{nginx}, \texttt{ubuntu} and
\texttt{bhuvnesh}. Not all of them are humans.

\begin{center}
\small
\begin{tabular}{@{}l l@{}}
\toprule
\rowcolor{tablehead}\bfseries User & \bfseries Purpose \\
\midrule
\texttt{root}     & Administrator \\
\texttt{mysql}    & MySQL service \\
\texttt{postgres} & PostgreSQL service \\
\texttt{nginx}    & Nginx server \\
\texttt{www-data} & Apache / Nginx web apps \\
\bottomrule
\end{tabular}
\end{center}

These are called \textbf{service accounts}.

\subsection{Process Ownership}

Every running process belongs to a user.

\begin{sh}
ps aux
\end{sh}

\begin{lstlisting}[language=]
USER      PID
root      123
postgres  456
nginx     789
\end{lstlisting}

PostgreSQL runs as \texttt{postgres} and Nginx runs as \texttt{nginx}. This improves
security.

\subsection{Useful Admin Commands}

\begin{cmdtable3}[3.4cm]
useradd & Create a user, with a home directory & sudo useradd -m john \nr
passwd & Set a user's password & sudo passwd john \nr
groupadd & Create a group & sudo groupadd developers \nr
usermod & Add a user to a group & sudo usermod -aG developers john \nr
chown & Change the owner of a file & sudo chown john file.txt \nr
chgrp & Change the group of a file & sudo chgrp developers file.txt \nr
userdel & Delete a user and their home directory & sudo userdel -r john \nr
groupdel & Delete a group & sudo groupdel developers \nr
whoami & Show the current user & whoami \nr
id & Show user and group IDs & id \\
\end{cmdtable3}

\subsection{How This Applies to Docker, Nginx and PostgreSQL}

When you run \lstinline|docker run nginx|, Nginx inside the container often runs as
\texttt{nginx}, not \texttt{root}. When PostgreSQL starts, \texttt{postgres} owns the
database files. When GitHub Actions deploys code, \texttt{ubuntu} or a dedicated deployment
user may own the files.

Understanding \textbf{users, groups, ownership and permissions} is one of the most important
Linux concepts, because it explains why commands sometimes fail with:

\begin{lstlisting}[language=]
Permission denied
\end{lstlisting}

Most of the time the current user is not the owner, is not in the required group, or lacks
the necessary permission bits (\texttt{r}, \texttt{w}, \texttt{x}).

\end{document}
